\daytitle{9}{机械能守恒与功能关系}

\begin{learningbox}{本讲义大致内容}
动能、重力势能、弹性势能统称为机械能。今天学习机械能守恒的条件，以及当机械能不守恒时如何用功能关系处理。能量方法是高中物理最重要的思想之一。
\end{learningbox}

\section*{一、学习目标}

\begin{enumerate}
  \item 掌握重力势能 $E_p = mgh$，理解重力做功与路径无关。
  \item 掌握弹性势能 $E_p = \frac12 kx^2$。
  \item 理解机械能守恒的条件：只有重力（或弹簧弹力）做功。
  \item 会用机械能守恒定律解决多过程问题。
  \item 掌握功能关系：除重力/弹力外其他力做功等于机械能变化量。
\end{enumerate}

\section*{二、初高中衔接}

初中已经学过"动能和势能可以相互转化"，如自由落体中重力势能转化为动能。但初中只讲定性，没有定量公式。

高中把这两个量定量化：
\begin{itemize}
  \item $E_k = \frac12 mv^2$（动能）——Day 08 已学。
  \item $E_p = mgh$（重力势能）——今天学。
  \item $E_p = \frac12 kx^2$（弹性势能）——今天学。
\end{itemize}

更关键的是，高中要判断\textbf{什么时候机械能守恒、什么时候不守恒}，以及不守恒时能量去哪了。

\section*{三、核心概念}

\begin{conceptbox}{重力势能}

\[
\boxed{E_p = mgh}
\]

\begin{itemize}
  \item $h$ 是物体相对参考平面的高度。
  \item 重力势能是相对的——参考平面不同，数值不同。
  \item 重力势能的变化是绝对的——$\Delta E_p = mg\Delta h$ 与参考面选择无关。
\end{itemize}

\textbf{重力做功与重力势能的关系：}
\[
W_G = -\Delta E_p = -(E_{p2} - E_{p1}) = E_{p1} - E_{p2}
\]

\begin{itemize}
  \item 重力做正功 $\to$ 重力势能减少。
  \item 重力做负功 $\to$ 重力势能增加。
  \item 重力做功与路径无关，只与初末高度差有关。
\end{itemize}

\end{conceptbox}

\begin{conceptbox}{弹性势能}

\[
\boxed{E_p = \frac12 k x^2}
\]

\begin{itemize}
  \item $k$：劲度系数（$\mathrm{N/m}$）。
  \item $x$：形变量（伸长或压缩的长度）。
  \item 弹力做功等于弹性势能变化的负值：$W_{\text{弹}} = -\Delta E_p$。
\end{itemize}

\textbf{注意：} 弹性势能永远是正值（因为平方）。弹簧被压缩和拉长相同距离时，弹性势能相同。

\end{conceptbox}

\begin{conceptbox}{机械能守恒定律}

\textbf{机械能：} 动能、重力势能、弹性势能的总和。
\[
E = E_k + E_p = \frac12 mv^2 + mgh + \frac12 kx^2
\]

\textbf{守恒条件：} 系统内只有重力（或弹簧弹力）做功，\textbf{且除重力、弹簧弹力外的其他力对系统做功为零}。注意：
\begin{itemize}
  \item "没有摩擦力"是一个常见判断捷径——物体在固定粗糙面上相对滑动时，滑动摩擦力做负功，机械能不守恒。但静摩擦力可以做正功（如传送带），此时机械能增加。
  \item "有外力"本身 \textbf{不等于} 机械能不守恒——关键是看这个外力是否做功。例如用手托着物体匀速上升，手对物体做正功，机械能不守恒；但若外力始终与位移垂直（如光滑面上的向心力），不做功，机械能仍守恒。
\end{itemize}

\textbf{守恒定律：}
\[
\boxed{E_1 = E_2}
\]
即：
\[
\frac12 mv_1^2 + mgh_1 + \frac12 kx_1^2 = \frac12 mv_2^2 + mgh_2 + \frac12 kx_2^2
\]

\begin{center}
% TikZ 图：机械能守恒（摆球）
% Day 09：机械能守恒与功能关系
\begin{tikzpicture}[scale=0.9,>=Latex]
  % 天花板
  \draw[thick] (-1.5,0) -- (1.5,0);
  % 悬点
  \filldraw[black] (0,0) circle (2pt);
  % 摆线 + 摆球（最高点）
  \draw[thick] (0,0) -- (-1.2,-2.5);
  \filldraw[black] (-1.2,-2.5) circle (3pt);
  \node[left] at (-1.2,-2.5) {$m$};
  \node[above left] at (-0.6,-1.25) {$l$};
  % 最低点
  \draw[dashed] (0,0) -- (0,-2.5);
  \filldraw[black] (0,-2.5) circle (2pt);
  \node[right] at (0,-2.5) {$v$};
  % 高度差
  \draw[<->] (-1.5,-2.5) -- (-1.5,-3.5);
  \node[left] at (-1.5,-3) {$h$};
  % 能量转换标注
  \node at (0,-3.8) {最高点：$E_p = mgh$，$E_k = 0$};
  \node at (0,-4.2) {最低点：$E_p = 0$，$E_k = \frac12 mv^2$};
  \node[red] at (0,-4.6) {$mgh = \frac12 mv^2$};
\end{tikzpicture}


\figcaption{摆球：最高点势能 $\to$ 最低点动能，机械能守恒}\label{fig:energy-conservation}
\end{center}

\textbf{常见判断：}
\begin{itemize}
  \item 自由落体、抛体运动（不计空气阻力）$\to$ 只有重力做功 $\to$ 机械能守恒。
  \item 光滑斜面上滑动 $\to$ 只有重力做功 $\to$ 机械能守恒。
  \item 粗糙面上滑动 $\to$ 摩擦力做负功 $\to$ 机械能不守恒，转化为内能。
  \item 外力做正功（如手拉物体上升）$\to$ 机械能增加。
  \item 外力不做功（如向心力）$\to$ 机械能守恒。
\end{itemize}

\end{conceptbox}

\begin{conceptbox}{功能关系}

当机械能不守恒时，其他力做的功等于机械能的变化量：
\[
\boxed{W_{\text{其他}} = \Delta E = E_2 - E_1}
\]

\begin{itemize}
  \item $W_{\text{其他}}$ 是除重力和弹簧弹力以外的力做的总功。
  \item 物体在固定粗糙面上相对滑动时，滑动摩擦力做负功 $\to$ 机械能减少，转化为内能。
  \item 外力做正功 $\to$ 机械能增加。

\end{itemize}

\textbf{常见模型：}
\begin{itemize}
  \item 粗糙斜面下滑：$mgh - \mu mg\cos\theta \cdot s = \frac12 mv^2$（动能定理）
  \item 或从功能关系看：$-\mu mg\cos\theta \cdot s = \Delta E$（机械能减少）
\end{itemize}

\end{conceptbox}

\section*{四、公式与适用条件}

\begin{center}
\begin{tabular}{lll}
\toprule
公式 & 意义 & 适用条件 \\
\midrule
$E_p = mgh$ & 重力势能 & $h$ 相对参考平面 \\
$E_p = \frac12 kx^2$ & 弹性势能 & 弹性限度内 \\
$E_1 = E_2$ & 机械能守恒 & 除重力/弹力外其他力做功为零 \\
$W_{\text{其他}} = \Delta E$ & 功能关系 & 选定系统，其他力总功 \\
\bottomrule
\end{tabular}
\end{center}

\section*{五、典型模型}

\subsection*{模型 1：自由落体/抛体}

只有重力做功，机械能守恒。$\frac12 mv_1^2 + mgh_1 = \frac12 mv_2^2 + mgh_2$。

\subsection*{模型 2：弹簧-物体系统}

物体与弹簧组成的系统，只有重力和弹力做功，机械能守恒。
\[
\frac12 mv^2 + mgh + \frac12 kx^2 = \text{常量}
\]

\subsection*{模型 3：粗糙斜面（功能关系）}

物体沿粗糙斜面下滑，摩擦力做负功，机械能减少：
\[
mgh - \mu mg\cos\theta \cdot s = \frac12 mv^2
\]

\section*{六、例题精讲}

\begin{examplebox}{机械能守恒——摆球}

\textbf{题意分析：} 摆长为 $l$ 的单摆，将摆球拉到与竖直方向成 $60^\circ$ 角的位置由静止释放，$g = 10\ \mathrm{m/s^2}$。求摆球经过最低点时的速度。

\textbf{物理模型：} 单摆，只有重力做功，机械能守恒。

\textbf{解题过程：}

取最低点为参考平面（$h=0$）。

初始高度：$h = l - l\cos 60^\circ = l(1 - 0.5) = 0.5l$

初态：$E_{p1} = mgh = 0.5 mgl$，$E_{k1} = 0$。

末态（最低点）：$E_{p2} = 0$，$E_{k2} = \frac12 mv^2$。

机械能守恒：
\[
0.5 mgl = \frac12 mv^2
\]
\[
v = \sqrt{gl}
\]

若 $l = 1.6\ \mathrm{m}$，则 $v = \sqrt{10 \times 1.6} = 4\ \mathrm{m/s}$。

\textbf{方法总结：}
\begin{enumerate}
  \item 机械能守恒的关键是选好参考平面（通常选最低点或最容易算的位置）。
  \item 摆球问题中高度差通过几何关系求出。
\end{enumerate}

\end{examplebox}

\begin{examplebox}{机械能守恒——弹簧+物体}

\textbf{题意分析：} 劲度系数 $k = 100\ \mathrm{N/m}$ 的轻弹簧竖直放置，下端固定。在弹簧上端放一个 $m = 0.5\ \mathrm{kg}$ 的物体（不拴接），从原长位置由静止释放，$g = 10\ \mathrm{m/s^2}$。求弹簧的最大压缩量。

\textbf{物理模型：} 物体+弹簧系统，只有重力和弹力做功，机械能守恒。

\textbf{解题过程：}

取弹簧原长位置为重力势能零点，此时弹簧无形变。

初态：$E_{k1}=0$，$E_{pG1}=0$，$E_{p\text{弹}1}=0$。

末态（最大压缩量 $x$ 时，速度为零）：$E_{k2}=0$，$E_{pG2} = -mgx$（物体下降了 $x$），$E_{p\text{弹}2} = \frac12 kx^2$。

机械能守恒：
\[
0 = -mgx + \frac12 kx^2
\]
\[
x( \frac12 kx - mg) = 0
\]
舍去 $x=0$，得：
\[
x = \frac{2mg}{k} = \frac{2 \times 0.5 \times 10}{100} = 0.1\ \mathrm{m} = 10\ \mathrm{cm}
\]

\textbf{方法总结：}
\begin{enumerate}
  \item 弹簧问题中重力势能和弹性势能都要考虑。
  \item 注意重力势能的正负——物体下降时 $h$ 为负，$E_p$ 为负。
\end{enumerate}

\end{examplebox}

\begin{examplebox}{功能关系——粗糙斜面}

\textbf{题意分析：} 质量为 $m = 2\ \mathrm{kg}$ 的物体从 $h = 3\ \mathrm{m}$ 高处沿粗糙斜面由静止滑下，斜面倾角 $\theta = 30^\circ$，$\mu = 0.3$，$g = 10\ \mathrm{m/s^2}$。求物体滑到底端时的速度。

\textbf{物理模型：} 粗糙斜面，摩擦力做负功，机械能不守恒。用动能定理或功能关系。

\textbf{解题过程：}

斜面长度 $s = \dfrac{h}{\sin\theta} = \dfrac{3}{0.5} = 6\ \mathrm{m}$。

重力做功：$W_G = mgh = 2 \times 10 \times 3 = 60\ \mathrm{J}$

摩擦力做功：$f = \mu mg\cos\theta = 0.3 \times 20 \times \frac{\sqrt{3}}{2} \approx 5.2\ \mathrm{N}$，$W_f = -f s \approx -31.2\ \mathrm{J}$

由动能定理：
\[
W_G + W_f = \frac12 mv^2
\]
\[
60 - 31.2 = \frac12 \times 2 \times v^2
\]
\[
v^2 = 28.8,\qquad v \approx 5.37\ \mathrm{m/s}
\]

从功能关系看：摩擦力做功 $-31.2\ \mathrm{J}$ 使机械能减少了 $31.2\ \mathrm{J}$，这些能量转化为内能。

\textbf{方法总结：}
\begin{enumerate}
  \item 机械能不守恒时，用动能定理（功能关系）仍然成立。
  \item 摩擦力做的功等于机械能的减少量，等于产生的热量。
\end{enumerate}

\end{examplebox}

\section*{七、变式训练}

\begin{enumerate}
  \item 质量为 $m = 0.2\ \mathrm{kg}$ 的小球从 $h = 0.8\ \mathrm{m}$ 高处自由下落到竖直放置的弹簧上，弹簧劲度系数 $k = 100\ \mathrm{N/m}$。求弹簧的最大压缩量。
  
  \item 一个物体从光滑斜面顶端由静止滑下，斜面高 $h = 1.8\ \mathrm{m}$，底端接粗糙水平面，$\mu = 0.3$，$g = 10\ \mathrm{m/s^2}$。求物体在水平面上滑行的距离。
  
  \item 质量为 $m$ 的小球用长为 $l$ 的细绳悬挂，将小球拉到水平位置由静止释放。求小球经过最低点时绳子的拉力。（提示：最低点时 $T - mg = m\frac{v^2}{l}$。）
\end{enumerate}

\section*{八、易错点}

\begin{mistakebox}
\begin{enumerate}
  \item \textbf{误以为只要物体运动就机械能守恒：} 判断条件不是"运动"而是"除重力/弹力外其他力做功是否为零"。摩擦力做负功时机械能减少，做正功时机械能增加。
  
  \item \textbf{重力势能的参考平面选错：} 同一个问题中，全程只能用同一个参考平面。换参考面会导致计算错误。
  
  \item \textbf{忽略弹性势能：} 有弹簧时，必须把弹性势能计入机械能的总和中。
  
  \item \textbf{把动能的 $v^2$ 写成 $v$：} $E_k = \frac12 mv^2$，不是 $\frac12 mv$。这是最常见的计算错误。
  
  \item \textbf{机械能守恒方程中符号搞反：} $\Delta E_k = -\Delta E_p$，动能的增加等于势能的减少，不是增加。
\end{enumerate}
\end{mistakebox}

\section*{九、竞赛入门补充}

\begin{competitionbox}
\textbf{系统的机械能守恒}

竞赛中经常通过绳、杆连接两个物体。虽然单个物体的机械能不守恒，但\textbf{系统的机械能守恒}（除重力/弹力外其他力对系统做功为零）。

\textbf{例题：} 质量 $m_A$ 和 $m_B$ 通过跨过定滑轮的轻绳连接，$A$ 在光滑水平面上，$B$ 从 $h$ 高处由静止释放。求 $B$ 落地时的速度。

\textbf{解：} 对 $A$ 和 $B$ 组成的系统，除重力外其他力对系统做功为零，系统机械能守恒：
\[
m_B g h = \frac12 m_A v^2 + \frac12 m_B v^2
\]
\[
v = \sqrt{\frac{2 m_B g h}{m_A + m_B}}
\]

注意两物体速度大小相等（绳不伸长），所以动能为两项之和。

\textbf{竞赛小贴士：} 连接体问题中，如果除重力和弹力外其他力对系统做功为零，优先用系统机械能守恒，比用牛顿定律简单得多。
\end{competitionbox}

\section*{十、今日小结}

\begin{enumerate}
  \item 重力势能 $E_p = mgh$，弹性势能 $E_p = \frac12 kx^2$。
  \item 机械能守恒条件：除重力/弹力外其他力对系统做功为零。
  \item 机械能守恒方程：$E_{k1} + E_{p1} = E_{k2} + E_{p2}$。
  \item 功能关系：$W_{\text{其他}} = \Delta E$，其他力做功等于机械能变化。
  \item 系统机械能守恒：连接体问题中优先考虑。
\end{enumerate}
